\begin{ouverture}[radioactivite]{10}{Transformations nucléaires}{Constitution et transformations}{Observer · Modéliser · Expérimenter · Expliquer}
\objectif{Reconnaître des isotopes à leur noyau.}
\objectif{Distinguer transformation physique, chimique et nucléaire.}
\objectif{Associer fusion et fission aux conversions d’énergie.}
\prerequis{Proportionnalité, conversions et lecture de graphiques.}
\end{ouverture}
\essentiel
\essSection{1}{Même élément, noyaux différents}
Des isotopes possèdent le même nombre de protons $Z$ et des nombres de neutrons différents. Ils appartiennent au même élément chimique, mais leurs nombres de masse $A$ diffèrent.
\begin{exemple}Les noyaux $\ce{^{12}_{6}C}$ et $\ce{^{14}_{6}C}$ possèdent tous deux six protons, mais respectivement six et huit neutrons. Ils sont isotopes. Tous les isotopes ne sont pas radioactifs.\end{exemple}
\essSection{2}{Le noyau se transforme}
Une transformation physique conserve les espèces ; une transformation chimique redistribue les atomes entre espèces sans modifier leurs noyaux ; une transformation nucléaire modifie des noyaux. Dans les équations nucléaires simples du livre, les sommes des nombres $A$ et des nombres $Z$ se conservent.
\schema[corrige]{nucleaire}
\begin{savoirfaire}[Lire une équation]Identifier les noyaux de départ et d'arrivée ; vérifier séparément les sommes de $A$ et de $Z$. Un neutron s'écrit $\ce{^{1}_{0}n}$. Une équation de fusion ou de fission peut comporter d'autres produits, comme des rayonnements, qui ne sont pas dessinés dans un modèle simplifié.\end{savoirfaire}
\essSection{3}{Fusion, fission et énergie}
La fusion assemble des noyaux légers ; la fission scinde un noyau lourd. Le Soleil convertit de l'énergie au cours de transformations nucléaires dont le bilan transforme de l'hydrogène en hélium. Dans une centrale actuelle à fission, la fission de noyaux lourds apporte de l'énergie thermique ; une chaîne de conversions permet ensuite de produire de l'électricité.
\begin{attention}La fusion nucléaire n'est pas la fusion d'un solide. Le même mot désigne ici une modification du noyau ; dans le changement d'état, les noyaux et les espèces restent inchangés.\end{attention}
On compare des bilans énergétiques fournis sans imposer ici la relation entre masse et énergie ni une loi de décroissance radioactive.
\essSection{4}{Expérimenter et modéliser}
\begin{experience}[Modéliser des bilans]Avec des jetons de deux couleurs représentant protons et neutrons, vérifier une équation fournie. Le modèle contrôle les nombres ; il ne simule ni les mécanismes, ni les conditions de déclenchement, ni la probabilité des réactions. Les études nucléaires du manuel sont documentaires, sans manipulation de source radioactive.\end{experience}
\exercices
\rubrique{Prendre confiance}
\begin{blocLivre}{ex-10-1}\begin{exercice}[1]{Reconnaître des isotopes}
On donne $\ce{^{16}_{8}O}$, $\ce{^{18}_{8}O}$ et $\ce{^{19}_{9}F}$.
\begin{questions}
\item Quels noyaux sont isotopes ?
\item Donner leurs nombres de neutrons.
\end{questions}
\end{exercice}\end{blocLivre}
\begin{blocLivre}{ex-10-2}\begin{exercice}[1]{Classer des transformations}
On observe une vaporisation, une combustion et une fission.
\begin{questions}
\item Associer physique, chimique et nucléaire.
\item Laquelle change la nature de certains noyaux ?
\end{questions}
\end{exercice}\end{blocLivre}
\begin{blocLivre}{ex-10-3}\begin{exercice}[1]{Un neutron}
Le neutron est noté $\ce{^{1}_{0}n}$.
\begin{questions}
\item Interpréter les nombres 1 et 0.
\item Quelle est sa charge électrique ?
\end{questions}
\end{exercice}\end{blocLivre}
\begin{blocLivre}{ex-10-4}\begin{exercice}[1]{La fusion de la glace}
Un texte affirme que faire fondre un glaçon est une réaction nucléaire de fusion.
\begin{questions}
\item Corriger cette phrase.
\item Les molécules d’eau sont-elles transformées en une autre espèce ?
\end{questions}
\end{exercice}\end{blocLivre}
\rubrique{S’entraîner}
\begin{blocLivre}{ex-10-5}\begin{exercice}[2]{Une fusion fournie}
$\ce{^{2}_{1}H + ^{3}_{1}H -> ^{4}_{2}He + ^{1}_{0}n}$.
\schema{nucleaire}
\begin{questions}
\item Vérifier les deux bilans.
\item Pourquoi s’agit-il d’une fusion nucléaire ?
\end{questions}
\end{exercice}\end{blocLivre}
\begin{blocLivre}{ex-10-6}\begin{exercice}[2]{Compléter une fission}
$\ce{^{235}_{92}U + ^{1}_{0}n -> ^{141}_{56}Ba + ^{92}_{36}Kr + x ^{1}_{0}n}$.
\begin{questions}
\item Déterminer $x$ par conservation de $A$.
\item Contrôler la somme de $Z$.
\end{questions}
\end{exercice}\end{blocLivre}
\begin{blocLivre}{ex-10-7}\begin{exercice}[2]{La source d’énergie du Soleil}
Le Soleil émet de l’énergie pendant des milliards d’années.
\begin{questions}
\item À quelle famille de transformations cette énergie est-elle reliée ?
\item Pourquoi une combustion chimique est-elle un autre modèle ?
\end{questions}
\end{exercice}\end{blocLivre}
\begin{blocLivre}{ex-10-8}\begin{exercice}[2]{Une chaîne de conversions}
Une centrale à fission fournit de l’électricité.
\begin{questions}
\item Compléter : nucléaire vers thermique vers mouvement vers électrique.
\item Pourquoi toute l’énergie ne devient-elle pas électrique ?
\end{questions}
\end{exercice}\end{blocLivre}
\rubrique{Raisonner et approfondir}
\begin{blocLivre}{ex-10-9}\begin{exercice}[3]{Un noyau inconnu}
$\ce{^{14}_{7}N + ^{4}_{2}He -> ^{17}_{8}O + X}$.
\begin{questions}
\item Déterminer $A$ et $Z$ de X.
\item Identifier X parmi proton, neutron et noyau d’hélium.
\end{questions}
\end{exercice}\end{blocLivre}
\begin{blocLivre}{ex-10-10}\begin{exercice}[3]{Un bilan énergétique fourni}
Une réaction donnée libère $3,2\times10^{-11}$ J par événement. Il se produit $2,0\times10^{18}$ événements.
\begin{questions}
\item Calculer l’énergie totale.
\item Quelle hypothèse permet la multiplication ?
\end{questions}
\end{exercice}\end{blocLivre}
\begin{blocLivre}{ex-10-11}\begin{exercice}[3]{Charge et nucléons}
Une équation proposée conserve $A$ mais transforme un noyau de $Z=10$ en deux noyaux de $Z=4$ et $Z=5$, sans autre particule chargée.
\begin{questions}
\item Quel bilan pose problème ?
\item Pourquoi un seul bilan correct ne suffit-il pas ?
\end{questions}
\end{exercice}\end{blocLivre}
\begin{blocLivre}{ex-10-12}\begin{exercice}[3]{Isotopes et ions}
On compare $\ce{^{23}_{11}Na}$ et $\ce{^{23}_{11}Na+}$.
\begin{questions}
\item S’agit-il de deux isotopes différents ?
\item Qu’est-ce qui diffère entre les deux entités ?
\end{questions}
\end{exercice}\end{blocLivre}

